\section{AlexNet 时刻：Software 2.0 从研究突破转为产业信号}

\subsection{为什么 2012 年成为拐点}

访谈重提了 2012 年 ImageNet 竞赛中 AlexNet 的压倒性结果。更重要的是 Jensen 描述了 NVIDIA 当时的内部情境：公司自己也在做计算机视觉，但传统方法进展受阻，AlexNet 的突破等于在内部痛点上“直接击中”。

\begin{figure}[H]
\centering
\includegraphics[width=0.88\textwidth]{frame-alexnet-slide.jpg}
\caption{访谈中展示 AlexNet 论文封面，强调其范式意义 \protect\footnotemark}
\end{figure}
\footnotetext{视频画面时间区间：00:11:20--00:11:30。}

\begin{knowledgebox}{从“编程”到“训练”的迁移}
Jensen 对该迁移的定义可以简化为：
\begin{itemize}
    \item Software 1.0：人写规则，机器执行规则。
    \item Software 2.0：人给数据，机器学习规则。
\end{itemize}

当问题规模和复杂性超过手写规则可维护边界时，2.0 路径会持续侵蚀 1.0 领域。这也是后续 LLM 爆发能迅速跨行业扩散的原因。
\end{knowledgebox}

\subsection{“Reasoned hope”与二次确认}

Jensen 用“reasoned hope”形容当年的判断：既不是无证据乐观，也不是保守等验证，而是在工程与市场约束内做高赔率投资。AlexNet 的出现给了 NVIDIA 二次确认：
\begin{itemize}
    \item 并行硬件路径方向正确。
    \item 训练型工作负载将成为长期主流。
    \item 软件生态应围绕学习系统而不是固定算法设计。
\end{itemize}

\begin{warningbox}{误区：AlexNet 只是“某个模型赢了比赛”}
把 AlexNet 看成一次 benchmark 胜利，会错过核心信息。真正被验证的是“数据 + 计算 + 可扩展模型”的共同规律，以及 GPU 在该规律里的基础设施地位。
\end{warningbox}

\subsection{本章小结}
AlexNet 对 NVIDIA 的意义不是“蹭到热点”，而是确认了其十年前下注的 CUDA 路线能够承载下一代计算范式。

